% =========================================================
% Proof: Preservation of Boundedness by Isometries
% Source: notes/isometries/notes-isometries.tex
% Difficulty: Easy
% =========================================================

\subsection*{Preservation of Boundedness}
\label{prf:isometry-preserves-bounded}

\begin{remark*}[Return]
$\leftarrow$ \hyperref[prop:isometry-preserves-bounded]{Back to Proposition in Notes}
\end{remark*}

\begin{proof}
\Claim If $f : X \to Y$ is an isometry and $A \subseteq X$, then $A$ is
bounded iff $f(A)$ is bounded.

\Given An isometry $f$ and $A \subseteq X$.

\Goal To prove $A$ bounded $\iff$ $f(A)$ bounded.

\Strategy A set is bounded iff its diameter is finite.
By Proposition~\ref{prop:isometry-preserves-diameter},
$\operatorname{diam}(f(A)) = \operatorname{diam}(A)$.
Finite iff finite.

% -------------------------------------------------------
% YOUR PROOF HERE
% -------------------------------------------------------

\end{proof}

\begin{remark*}[Proof shape]
Immediate from diameter preservation: bounded = finite diameter.
\end{remark*}

\begin{remark*}[Dependencies]
\begin{itemize}
  \item Proposition~\ref{prop:isometry-preserves-diameter}: $\operatorname{diam}(f(A)) = \operatorname{diam}(A)$.
  \item Definition~\ref{def:bounded-set}: bounded iff finite diameter.
\end{itemize}
\end{remark*}
